\begin{musica}{Fado do Estudante}{Vasco Santana}{}
\ac{C}{Tri}ste sina ver-me as\ac{G7}{sim}
Que sorte vil degra\ac{C}{dan}te.
Ai, que saudade eu sinto em \ac{G7}{mim}
Do meu viver de estu\ac{C}{dan}te.
\ac{C7}{De}sse fu\ac{F}{gaz} tempo de amor
Que d'um ra\ac{C}{paz} era o melhor,
Era um au\ac{G7}{daz} conquistador das rapa\ac{C}{ri}gas.
\ac{C7}{De} capa ao \ac{F}{ar} cabeça ao léu,
Só para a\ac{C}{mar} vivia eu
Sem me ra\ac{G7}{lar} e tudo mais eram can\ac{C}{ti}gas.

Nenhuma delas me prendeu,
Deixá-las eu era canja;
Até ao dia em que apareceu
Essa traidora de franja.

Sempre a tenir sem um tostão,
Batina a abrir por um rasgão,
Botas a rir e um bengalão, e ar descarado.
A vadiar com outros mais
Para ir dançar para os arraias
P’ra namorar, beber, folgar, cantar o fado.

Recordo agora com saudade
Os calhamaços que eu lia,
Os professores da faculdade
E a mesa de anatomia.

Evoco em mim recordações
Que não têm fim dessas lições
Frente ao jardim do velho campo de santana;
Aulas que dava e sem estudar
Só ainda estava nessa classe
A que eu faltava sete dias por semana.

O fado é toda a minha fé,
Embala, encanta, inebria.
Pois chega a ser bonito até
Na rádio-telefonia.

Quando é tocado com amor,
Bem afinado e a rigor
É belo fado, ningém há que lhe resista;
É a canção mais popular,
Tem emoção, faz-nos vibrar,
E eis a razão de eu ser doutor e ser.

\posicoes{C: x32010, C7: x32310, F: 133211, G7: 320001}
\end{musica}
